\documentclass[10pt,a4paper]{amsart}
\usepackage[margin=19mm]{geometry}
\usepackage{amsmath,amssymb,mathtools}
\usepackage[scr=rsfs,cal=euler]{mathalfa}
\usepackage{xfrac}
\usepackage[unicode,hidelinks,bookmarks=false]{hyperref}
\newcommand{\ie}{i.e.\ }
\newcommand{\cf}{cf.\ }
\newcommand{\resp}{respectively\ }
\newcommand{\Subsection}{\subsection}
\providecommand{\nobreakdash}{\nobreak\hbox{-}}
\setlength{\emergencystretch}{2em}
\multlinegap0pt
\usepackage{amssymb}
\usepackage{xcolor}
\newtheorem{definition}{Definition}
\newtheorem{lemma}{Lemma}
\newtheorem{theorem}{Theorem}
\title{A pathological set regarding the propagation of almost sure properties of Gaussian measures}
\author{Pablo Merino}
\date{Source: arXiv:2605.31359v1 (CC BY 4.0)}
\begin{document}
\maketitle

\noindent\emph{Source and licence.} A pathological set regarding the propagation of almost sure properties of Gaussian measures. Version arXiv:2605.31359v1 (CC BY 4.0), distributed by arXiv under the Creative Commons Attribution 4.0 International licence, http://creativecommons.org/licenses/by/4.0/, as recorded in the arXiv licence metadata for this exact version and shown on the abstract page https://arxiv.org/abs/2605.31359v1. Full licence notice: \texttt{licenses/arxiv-2605.31359-CC-BY-4.0.txt}.\par\smallskip

\noindent\textit{Editorial scope.} Counted unit: the article's theorem thm:flambda-tnotzero (source0.tex lines 458--491). The article's complete original proof of it is delivered with it (source0.tex lines 493--578, one \texttt{proof} environment of 86 lines). No mathematics is written, repaired or re-derived: the statement and the proof are the article's own lines, in the article's own notation, and no retained line is substituted or re-flowed.  Retained uncounted as the source-local context the proof needs: lines 292--294; lines 335--338; lines 384--394; lines 404--413; lines 997--1002; lines 1005--1007.  Nothing was omitted from any retained span, so the package carries no omission record.  No retained line contains a citation, so the wrapper carries no bibliography and no bibliography entry.  No retained line contains a figure environment, an image-inclusion command or drawing code, so no image is delivered and the package declares no asset dependency.  Editorial formatting only, in addition: an AMS \texttt{amsart} preamble that restates the article's own declarations and macros used by the retained lines, namely the environments \texttt{definition}, \texttt{lemma}, \texttt{theorem} on separate counters, together with the standard packages those lines need.  The retained environments print in this document as Definition 1, Lemma 1, Theorem 1 in order of appearance, whereas the article prints the same items with the numbering of its own full document.  The complete native source of the article as distributed in the arXiv e-print for this exact version is bundled unchanged as \texttt{sources/201717/source0.tex}.
\begin{definition}\label{def:g-lambda}
Let $d \in \mathbb{N}$ and $\lambda \in (0,d)$. Define $F_{\lambda} = \mathcal{F}^{-1}[g_{\lambda}]$, where $g_{\lambda}(\xi) = \langle \xi \rangle^{-\sigma+1} h(\xi)|\xi|^{-\lambda}$.
\end{definition}

\begin{align}
			F_{\lambda} = \sum_{j \geq -1} P_{j}F_{\lambda},
			\label{eq:lp-decomp}
		\end{align}

\begin{lemma}\label{lemma:ftransformradial_bessel}
		Let $p \in [1,\infty)$. Suppose $f$ is a radial function in  $L^p(\mathbb{R}^d)$, $d \geq 2$; thus, $f(x) = f_0(|x|)$ for a.e. $x \in \mathbb{R}^d$. Moreover, assume that the integral
\begin{align}\label{eq:convcond-radialexp}
\int_0^{\infty} f_0(\rho) J_{\frac{d-2}{2}}(\rho|\xi|) \rho^{d/2} d\rho
\end{align}		
converges in a.e. $\xi \in \mathbb{R}^d$. Then the Fourier transform $\mathcal{F}[f]$ is given also by a radial function, and it has the form
		\begin{align*}
			\mathcal{F}[f](\xi) = c(d) |\xi|^{-\frac{d-2}{2}} \int_0^{\infty} f_0(\rho) J_{\frac{d-2}{2}}(\rho|\xi|) \rho^{d/2} d\rho \quad \text{for a.e. } \xi \in \mathbb{R}^d,
		\end{align*}
		where $c(d)$ is some constant depending only on the dimension.
	\end{lemma}

	\begin{lemma}\label{lemma:invft-xalpha-nu}
	Let $\alpha \in \mathbb{C}$ such that $\text{Re}(\alpha) \notin -\mathbb{N}$. Then, the tempered distributions $x^{\alpha} \mathbb{I}_{(1,\infty)}(x)$ admit the inverse distributional Fourier transform
	\begin{align*}
	&\mathcal{F}^{-1}[x^{\alpha} \mathbb{I}_{(1,\infty)}(x)](\theta)= J(\alpha) \{ e^{i\frac{\pi}{2}(\alpha+1)} \theta_+^{-\alpha-1} + e^{-i\frac{\pi}{2}(\alpha+1)} \theta_-^{-\alpha-1} \}
	\end{align*}
	where $J : \mathbb{C} \rightarrow \mathbb{C}$ is a function such that
	\begin{align*}
	|J(\alpha)| \leq \Gamma(\text{Re}(\alpha)+1) + \frac{\pi}{2}.
	\end{align*}
	\end{lemma}

	\begin{definition}\label{lemma:bessel_kreal}
		Let $k > -1/2$ and $t > 0$. Then
		\begin{align}\label{eq:bessel2}
			J_k(t) = \frac{(t/2)^k}{\Gamma(2k+1)\Gamma(1/2)} \int_{-1}^1 e^{its} (1-s^2)^{(2k-1)/2}ds.
		\end{align}
	\end{definition}

		\begin{align}\label{eq:asympBessel}
			J_k(t) = t^{-1/2} e^{it} \sum_{j=0}^N a_j t^{-j} + t^{-1/2} e^{-it} \sum_{j=0}^N b_j t^{-j} + R_k(t)
		\end{align}

\begin{theorem}\label{thm:flambda-tnotzero}
Let $d>1$, $t \in \mathbb{R} \setminus \{0\}$, $\lambda \in (\frac{d-1}{2},d)$ such that $\frac{d-1}{2} - \lambda \notin -\mathbb{N}$, and $J$ the function given in Lemma \ref{lemma:invft-xalpha-nu}. Then the following holds:
\begin{itemize}
\item For a.e. $x \in \mathbb{R}^d \setminus \{0\}$ such that $|x| \neq t$ and $-|x| \neq t$:
\begin{align*}
\langle &D \rangle^{\sigma-1} e^{it|D|} F_{\lambda}(x) \nonumber \\
&= J\left(\frac{d-1}{2}-\lambda\right)  |x|^{-\frac{d-1}{2}} \left\{ c_+ [ e^{i\frac{\pi}{2}(\frac{d+1}{2} - \lambda)} (|x| + t)_+^{\lambda-\frac{d+1}{2}}\right. \\
&\hspace{48mm} + e^{-i\frac{\pi}{2}(\frac{d+1}{2} - \lambda)} \left(|x| + t\right)_-^{\lambda-\frac{d+1}{2}}] \nonumber \\
&\left. + c_-  [ e^{i\frac{\pi}{2}(\frac{d+1}{2} - \lambda)}(-|x| + t)_+^{\lambda-\frac{d+1}{2}} + e^{-i\frac{\pi}{2}(\frac{d+1}{2} - \lambda)} \left( -|x| + t\right)_-^{\lambda-\frac{d+1}{2}} ] \right\} \nonumber  \\
&+ c_{\lambda}(x) + D_{\lambda}(x),
\end{align*}
for $c_+$ and $c_-$ constants, $D_{\lambda} \in C^{\infty}(\mathbb{R}^d)$ and 
\begin{align*}
|c_{\lambda}(x)| \lesssim \min\{|x|^{-\frac{d+1}{2}} , |x|^{-\frac{d-1}{2}}\}
\end{align*}
if $\lambda \in (\frac{d+1}{2},d)$, and 
\begin{align*}
|c_{\lambda}(x)| \lesssim |x|^{-\frac{d+1}{2}}
\end{align*}
if $\lambda \in (\frac{d-1}{2},d)$.
%; and if $x = 0$,
%\begin{align*}
%&\langle D \rangle^{\sigma} e^{\pm i t |D|} F_{\lambda}(0) = w_{d-1}\{ \Gamma(d - \lambda) + J(n-\lambda-1,\nu) \} \nonumber \\
%& \times \{ c_+ e^{i\frac{\pi}{2}(n-\lambda)} t_+^{-n + \lambda} + c_- e^{-i\frac{\pi}{2}(n-\lambda)} t_-^{-n + \lambda}\},
%\end{align*}
%for $w_{d-1}$ the surface of the $(d-1)$-dimensional sphere.

\item  Moreover, for any $A > 1$, $x \in \mathbb{R}^d \setminus B(0,A|t|)$ and $N > d - \lambda$, we have that
\begin{align*}
|\langle D \rangle^{\sigma-1} e^{it|D|} F_{\lambda}(x)| \lesssim_{d,\lambda} |x|^{-N}.
\end{align*}

\end{itemize}
\end{theorem}

\begin{proof}
Let $A > 1$ and $x \in \overline{B}(0,A|t|)^c$. We can repeat the procedure from the Littlewood-Paley decomposition in \eqref{eq:lp-decomp} with a slight modification. Indeed, maintaining the same notation,
\begin{align*}
&|\langle D \rangle^{\sigma-1} e^{it|D|} F_{\lambda}(x)| \leq \sum_{j \geq -1} \left| \int_{B_j} e^{i \xi x + t |\xi|}(\varphi(\xi/2^j) |\xi|^{-\lambda} h(\xi)) d\xi \right| \nonumber \\
&= \sum_{j \geq 1} 2^{j(d-\lambda)} \left| \int_{B_0} e^{i 2^j |x| (\frac{\eta x + t |\eta|}{|x|})}(\varphi(\eta) |\eta|^{-\lambda} ) d\eta \right| \nonumber \\
&+ \sum_{j=-1}^0 2^{j(d-\lambda)} \left| \int_{B_0} e^{i 2^j |x| (\frac{\eta x + t |\eta|}{|x|})}(\varphi(\eta) |\eta|^{-\lambda} h(2^j \eta)) d\eta \right|
\end{align*}
The phase $\phi(\xi) = \frac{\eta x + t |\eta|}{|x|}$ does not have critical points for $|x| > |t|$. Indeed,
\begin{align*}
\nabla_{\eta} \phi(\eta) = \frac{1}{|x|} \left( x + \frac{t \eta}{|\eta|} \right),
\end{align*}
so $\nabla_{\eta} \phi(\eta) = 0$ would imply $|x| = |t|$. Moreover, since $|x| > A |t|$, we have that for any $\eta \in B_0$
\begin{align*}
|\nabla_{\eta} \phi(\eta)| \geq 1 - \frac{|t|}{|x|} > 1 - \frac{1}{A}.
\end{align*}
Moreover, $\varphi(\eta) |\eta|^{-\lambda} h(2^j \eta)$ is smooth and compactly supported in $B_0$ for any $j \geq -1$ (notice that $h(2^j \cdot) = 1$ for $j \geq 1$).  Thus we can do the standard integration by parts to obtain
\begin{align*}
|\langle D \rangle^{\sigma-1} e^{it|D|} F_{\lambda}(x)| \lesssim_{d,\lambda,\varphi} |x|^{-N} \left( \sum_{j \geq -1} 2^{j(d-\lambda-N)} \right),
\end{align*}
where $N > d - \lambda$ suffices for the convergence of the right-hand side. %Note that, in this procedure, we did not use the presence of $e^{-\gamma |\xi|}$ in the integrand.

Let $x \in \mathbb{R}^d \setminus \{0\}$ such that $|x| \neq t$ and $-|x| \neq t$. Since $\lambda < n$, there exists some $q > 1$ such that $\lambda > \frac{d}{q}$. Recalling Definition \ref{def:g-lambda}, this implies that 
\begin{align*}
\langle \xi \rangle^{\sigma-1} e^{it|\xi|} g_{\lambda}(\xi) \in L^q(\mathbb{R}^d).
\end{align*}
Then, for almost every $x \in \mathbb{R}^d$, from Lemma \ref{lemma:ftransformradial_bessel} we have that
\begin{align*}
&\langle D \rangle^{\sigma-1} e^{it|D|} F_{\lambda}(x) \nonumber \\
&=c(d) |x|^{-\frac{d-2}{2}} \int_1^{\infty}  e^{it\rho}  \rho^{-\lambda} J_{\frac{d-2}{2}}(\rho|x|) \rho^{d/2} d\rho + \int_{|\xi| \leq 1} e^{ix\xi + it|\xi|} |\xi|^{-\lambda} h(\xi) d\xi \nonumber \\
&=(i) + (ii).  
\end{align*}
The object $(ii)$ is in $C^{\infty}(\mathbb{R}^d)$ for each $t \neq 0$ and $\lambda < n$. We focus on $(i)$. From Bessel functions asymptotic expansion \eqref{eq:asympBessel}, we have that
\begin{align*}
J_{\frac{d-2}{2}}(s)  - \left( c_+ s^{-1/2} e^{is} + c_- s^{-1/2} e^{-is} \right) = R_{\frac{d-2}{2}}(s)
\end{align*}
for $R_{\frac{d-2}{2}}(s) = \mathcal{O}(s^{-3/2})$ when $s \rightarrow \infty$, and 
$c_+, c_- \in \mathbb{C}$ constants. This expansion and the Poisson representation of Bessel functions given in Lemma \ref{lemma:bessel_kreal} implies that, for any $s \in \mathbb{R} \setminus \{0\}$,
\begin{align}\label{eq:ourbessel-asymp}
|s^{\frac{1}{2}}R_{\frac{d-2}{2}}(s)| \lesssim \min\{ 1 , s^{-1} \}.
\end{align}
Observe that $\frac{d-2}{2} > -\frac{1}{2}$ if and only if $d > 1$, i.e. $d>1$ is required for a well definition of the Bessel functions as we have defined them. We write
\begin{align*}
&(i) = |x|^{-\frac{d-2}{2}} \left( c_+ \int_1^{\infty}  e^{i\rho|x| + it\rho} \rho^{-\lambda} (\rho|x|)^{-1/2} \rho^{d/2} d\rho \right.   \nonumber \\
&\left. + c_- \int_1^{\infty} e^{-i\rho|x| + it\rho} \rho^{-\lambda} (\rho|x|)^{-1/2} \rho^{d/2} d\rho \right) + |x|^{-\frac{d-2}{2}} \int_1^{\infty} e^{it\rho} \rho^{-\lambda} R_{\frac{d-2}{2}}(\rho|x|) \rho^{d/2} d\rho \nonumber \\
&= |x|^{-\frac{d-1}{2}} \left(c_+ \int_{1}^{\infty}  e^{i \rho ( t + |x|)}  \rho^{\frac{d-1}{2} - \lambda} d\rho + c_- \int_{1}^{\infty}  e^{i \rho( t - |x|)}  \rho^{\frac{d-1}{2} - \lambda} d\rho \right) \nonumber \\
&+ |x|^{-\frac{d-2}{2}} \int_{1}^{\infty} e^{it\rho} R_{\frac{d-2}{2}}(\rho |x|) \rho^{\frac{d}{2} -\lambda} d\rho   = |x|^{-\frac{d-1}{2}}(*) + |x|^{-\frac{d-2}{2}}(**).
\end{align*}
Assume that $\frac{d-1}{2} - \lambda \notin -\mathbb{N}$. Thus, for $(*)$ we can use Lemma \ref{lemma:invft-xalpha-nu} with $\alpha = \frac{d-1}{2} - \lambda$, and $\theta = t + |x|$ for the first term and $\theta = t - |x|$ for the second one:
\begin{align*}
(*) &= J\left(\frac{d-1}{2}-\lambda\right) \left\{c_+ [ e^{i\frac{\pi}{2}(\frac{d+1}{2} - \lambda)} (|x| + t)_+^{\lambda-\frac{d+1}{2}}\right. \\
&\hspace{48mm} + e^{-i\frac{\pi}{2}(\frac{d+1}{2} - \lambda)} \left(|x| + t\right)_-^{\lambda-\frac{d+1}{2}}] \nonumber \\
&\hspace{48mm} + c_-  [ e^{i\frac{\pi}{2}(\frac{d+1}{2} - \lambda)}(-|x| + t)_+^{\lambda-\frac{d+1}{2}} \\
&\left. \hspace{48mm}+ e^{-i\frac{\pi}{2}(\frac{d+1}{2} - \lambda)} \left( -|x| + t\right)_-^{\lambda-\frac{d+1}{2}} ] \right\}
\end{align*}
With respect to $(**)$, we consider the inequality \eqref{eq:ourbessel-asymp} to obtain that
\begin{align*}
&|x|^{-\frac{d-2}{2}}\left|  (**) \right| \lesssim |x|^{-\frac{d-1}{2}} \int_{1}^{\infty}\rho^{\frac{d-1}{2} - \lambda} \min\{ 1 , (\rho|x|)^{-1} \} d\rho \nonumber.
\end{align*}
If we bound the integrand by $\min\{ 1 , (\rho|x|)^{-1} \} \leq (\rho|x|)^{-1}$ we obtain
\begin{align*}
|x|^{-\frac{d-2}{2}}\left|  (**) \right| \lesssim \left( \int_{1}^{\infty} \rho^{\frac{d-3}{2} - \lambda} d\rho \right) |x|^{-\frac{d+1}{2}},
\end{align*}
where the integral in the right hand side converges if $\lambda \in (\frac{d-1}{2},d)$. On the other hand, if we bound the integrand by $\min\{ 1 , (\rho|x|)^{-1} \} \leq 1$, then we obtain
\begin{align*}
|x|^{-\frac{d-2}{2}}\left|  (**) \right| \lesssim \left( \int_1^{\infty} \rho^{\frac{d-1}{2} - \lambda} d\rho \right) |x|^{-\frac{d-1}{2}}.
\end{align*}
Now, the integral in the right hand side of the last inequality converges if $\lambda \in (\frac{d+1}{2},d)$. Denote $|x|^{-\frac{d-2}{2}} (**) = c_{\lambda}(x)$, so that 
\begin{align*}
|c_{\lambda}(x)| \lesssim \min\{|x|^{-\frac{d+1}{2}} , |x|^{-\frac{d-1}{2}}\}
\end{align*}
if $\lambda \in (\frac{d+1}{2},d)$, and 
\begin{align*}
|c_{\lambda}(x)| \lesssim |x|^{-\frac{d+1}{2}}
\end{align*}
if $\lambda \in (\frac{d-1}{2},d)$. This concludes the proof.
%Regarding $x = 0$, observe that
%\begin{align*}
%&\langle D \rangle^{\sigma} e^{\pm i t |D|} F_{\lambda}(0) = \int_{B(0,\nu)^c} e^{\pm i t |\xi|} |\xi|^{-\lambda} d\xi = w_{d-1} \int_{\mathbb{R}} \mathbb{I}_{(\nu,\infty)} \rho^{d-1-\lambda} e^{\pm i t \rho}  d\rho \nonumber \\
%&=w_{d-1} \mathcal{F}[\mathbb{I}_{(\nu,\infty)} \rho^{d-1-\lambda}](\pm t).
%\end{align*}
%By result \eqref{lemma:invft-xalpha-nu},
%\begin{align*}
%&\langle D \rangle^{\sigma} e^{\pm i t |D|} F_{\lambda}(0) =  w_{d-1}\{ \Gamma(d - \lambda) + J(n-\lambda-1,\nu) \}  \nonumber \\
%&\times \{ c_+ e^{i \frac{\pi}{2}(n-\lambda)} t_+^{-n + \lambda} +  c_- e^{-i \frac{\pi}{2}(n-\lambda)} t_-^{-n + \lambda} \}.
%\end{align*}
\end{proof}

\end{document}
